%%%%%%%%%%%%%%%%%%%%%%%%%%%%%%%%%%%%%%%%%%%%%%%%%%%%%%%%%%%%%%%%%%%%%%%%%%%%%%%%
%  Policy Evaluation Lab -- Lecture 8: the macro toolkit, from GDP to IS-LM
%  Monday 16 November 2026, first lecture of Part II.
%  Built from Blanchard, Macroeconomics, 7th ed. (Pearson, 2017, Global
%  Edition): sections 2-1, 3-1 to 3-5, 4-1, 4-2, 4-4, 5-1 to 5-5, 9-1.
%  The book is copyrighted: wording, numerical examples and diagrams on these
%  slides are our own; notation and structure follow the book. Do not add
%  the PDF or its figures to the repository.
%  Checked against the book:
%   - three definitions of GDP; nominal vs real; growth (Y_t - Y_{t-1})/Y_{t-1}
%   - Table 3-1, US GDP 2014: C 68.3%, I 16.0%, G 18.1%, X 13.5%, IM 16.6%,
%     inventory investment 0.4% (Survey of Current Business, July 2015)
%   - G excludes transfers and interest payments
%   - Z = C + I + G; C = c0 + c1 (Y - T); Y = [c0 + I + G - c1 T]/(1 - c1)
%   - S = Y - T - C; I = S + (T - G); paradox of saving
%   - "Is the government omnipotent?": lags, other responses, expectations,
%     side effects (inflation), deficits and debt
%   - M^d = $Y L(i); M = $Y L(i); open market operations; P_B = 100/(1+i)
%   - central banks choose i; zero lower bound / liquidity trap
%   - I = I(Y, i); IS: Y = C(Y - T) + I(Y, i) + G; LM: M/P = Y L(i), and in
%     the 7th edition a horizontal LM at the rate chosen by the central bank
%   - three-step method; fiscal contraction; monetary expansion; policy mix;
%     fiscal consolidation + monetary expansion (Clinton and Greenspan, 1990s)
%   - Christiano, Eichenbaum and Evans (1996), as reported in Figure 5-10:
%     +1 point federal funds rate; retail sales -0.9% after 5 quarters;
%     output -0.7% after 8 quarters; price level nearly unchanged for about
%     6 quarters
%   - preface: AD-AS replaced by IS-LM-PC; eq. 9.4: pi - pi(-1) = (a/L)(Y - Yn)
%  Not used: the book's box on the German and French recessions (not checked).
%%%%%%%%%%%%%%%%%%%%%%%%%%%%%%%%%%%%%%%%%%%%%%%%%%%%%%%%%%%%%%%%%%%%%%%%%%%%%%%%

\documentclass[aspectratio=169,xcolor=table]{beamer}
\input{preamble-pel}
\usetikzlibrary{positioning,arrows.meta,calc}
\newcolumntype{L}[1]{>{\raggedright\arraybackslash}p{#1}}

\tikzset{
  ax/.style={-{Stealth[length=2.2mm]}, thick, pelText},
  curveA/.style={very thick, pelAccent},
  curveB/.style={very thick, pelSlate},
  shifted/.style={very thick, pelAccent, dashed},
  guide/.style={densely dotted, pelMuted},
  lbl/.style={font=\footnotesize},
}

\title[Policy Evaluation Lab]{Policy Evaluation Lab}
\subtitle[Lecture 8]{Lecture 8: The macro toolkit, from GDP to IS-LM}
\date[\pelYear]{November, 2026}

\begin{document}

\frame{\titlepage}

\begin{frame}{Part II starts today}
\small
\bi
  \item Part I: did \emph{this} policy work for \emph{these} firms, students, teachers?
  \item Part II: macroeconomic policy. What can fiscal and monetary policy do, who decides, under which rules
  \item Before we talk about the EU fiscal rules and the ECB, we need a common language:
  \bi
    \item what we measure (GDP and its components)
    \item how demand determines output in the short run
    \item how the central bank sets the interest rate
    \item how the two fit together: the \textbf{IS-LM} model
  \ei
\ei
\vspace{0.2cm}
\begin{block}{}
\centering One lecture, one model. A review, not a full course in macroeconomics
\end{block}
\pelTableNote{Notation and structure follow Blanchard, \textit{Macroeconomics}, 7th ed., chapters 2 to 5. Examples and diagrams are our own.}
\end{frame}

\begin{frame}{Today}
\be
  \item GDP: what it measures
  \item The goods market: demand, equilibrium output, the multiplier
  \item Investment equals saving
  \item Financial markets: money demand, money supply, the interest rate
  \item IS-LM: fiscal policy, monetary policy, the policy mix
  \item What the model leaves out, and where we go next
\ee
\end{frame}

%=============================================================================
\section{GDP}
%=============================================================================

\begin{frame}{GDP: three ways to look at the same number}
\small
A toy economy: a farmer grows wheat, a miller makes flour, a baker sells bread
\vspace{0.1cm}

\centering\footnotesize
\begin{tabular}{lrrrr}
\toprule
 & \textbf{Sales} & \textbf{Inputs bought} & \textbf{Value added} & \textbf{of which: wages / profits}\\
\midrule
Farmer (wheat) & 30 & 0 & 30 & 20 / 10 \\
Miller (flour) & 50 & 30 & 20 & 12 / 8 \\
Baker (bread)  & 90 & 50 & 40 & 28 / 12 \\
\midrule
\textbf{Total} & 170 & & \textbf{90} & 60 / 30 \\
\bottomrule
\end{tabular}
\vspace{0.15cm}

\raggedright\small
\be
  \item Value of \textbf{final} goods: the bread, 90 (wheat and flour are intermediate goods)
  \item Sum of \textbf{value added}: $30 + 20 + 40 = 90$
  \item Sum of \textbf{incomes}: wages 60 + profits 30 = 90
\ee
Adding up all sales (170) would count the wheat three times
\pelTableNote{Numbers invented for the example. Production and income are the same thing seen from two sides: this is why we use one symbol, $Y$, for both.}
\end{frame}

\begin{frame}{Nominal and real GDP}
\small
\bi
  \item \textbf{Nominal GDP} ($\$Y_t$): quantities $\times$ current prices. Grows if we produce more \emph{or} if prices rise
  \item \textbf{Real GDP} ($Y_t$): quantities $\times$ the prices of a base year
  \item GDP growth: $g_t = (Y_t - Y_{t-1})/Y_{t-1}$, computed on real GDP
\ei
\vspace{0.1cm}
\centering\footnotesize
\begin{tabular}{lrrrr}
\toprule
\textbf{Year} & \textbf{Loaves} & \textbf{Price} & \textbf{Nominal GDP} & \textbf{Real GDP (year 1 prices)}\\
\midrule
1 & 100 & 2.00 & 200 & 200 \\
2 & 105 & 2.10 & 220.5 & 210 \\
\bottomrule
\end{tabular}
\vspace{0.1cm}

\raggedright\small
Nominal growth 10.25\%, real growth 5\%: the rest is prices
\bi
  \item Positive growth: \textbf{expansion}. Negative growth: \textbf{recession}
  \item Real GDP per person: a (rough) measure of the standard of living
\ei
\end{frame}

\begin{frame}{Who buys what we produce}
\begin{columns}[T]
\begin{column}{0.47\textwidth}
\footnotesize
\[
Y = C + I + G + X - IM + \text{inventory inv.}
\]
\bi
  \item $C$: consumption, by far the largest
  \item $I$: \textbf{investment} = new machines, plants, houses. Not buying shares: that is \emph{financial} investment
  \item $G$: government purchases of goods and services. \textbf{Not} pensions, transfers, interest on debt
  \item $X - IM$: net exports
  \item Inventory investment: production minus sales
\ei
\end{column}
\begin{column}{0.5\textwidth}
\centering\footnotesize
\begin{tabular}{lr}
\toprule
\textbf{United States, 2014} & \textbf{\% of GDP}\\
\midrule
Consumption ($C$) & 68.3 \\
Investment ($I$) & 16.0 \\
Government spending ($G$) & 18.1 \\
Exports ($X$) & 13.5 \\
Imports ($IM$) & $-16.6$ \\
Inventory investment & 0.4 \\
\bottomrule
\end{tabular}
\pelTableNote{Source: Blanchard, Table 3-1, from the Survey of Current Business, July 2015.}
\vspace{0.1cm}
\raggedright\small
\footnotesize With transfers and interest, US public spending was about 33\% of GDP
\end{column}
\end{columns}
\end{frame}

%=============================================================================
\section{The goods market}
%=============================================================================

\begin{frame}{The demand for goods}
\footnotesize
Simplifications (for now): one good; firms supply whatever is demanded at a given price level $P$; a closed economy ($X = IM = 0$)
\[
Z \equiv C + I + G
\]
\bi
  \item Consumption depends on \textbf{disposable income} $Y_D = Y - T$ ($T$: taxes minus transfers)
  \[
  C = c_0 + c_1 (Y - T), \qquad 0 < c_1 < 1
  \]
  \bi
    \item $c_1$: propensity to consume. Out of one extra euro, how many cents are spent?
    \item $c_0$: consumption not explained by income (confidence, access to credit)
  \ei
  \item Investment $\bar I$: given, for now
  \item $G$ and $T$: chosen by the government. This is \textbf{fiscal policy}
\ei
\pelTableNote{Three kinds of equations: identities ($Z \equiv C+I+G$, $Y_D \equiv Y-T$), behavioural equations ($C = c_0 + c_1 Y_D$), equilibrium conditions ($Y = Z$, next slide).}
\end{frame}

\begin{frame}{Equilibrium output}
\small
Equilibrium: production equals demand, $Y = Z$
\[
Y = c_0 + c_1 (Y - T) + \bar I + G
\]
Solve for $Y$:
\[
Y = \underbrace{\frac{1}{1-c_1}}_{\text{multiplier}} \; \underbrace{\left[\,c_0 + \bar I + G - c_1 T\,\right]}_{\text{autonomous spending}}
\]
\bi
  \item Autonomous spending: the part of demand that does not depend on income
  \item Since $0 < c_1 < 1$, the multiplier is larger than 1
  \item With $c_1 = 0.75$: multiplier $= 1/0.25 = 4$. One more euro of $G$ (or of $c_0$, or of $\bar I$) raises output by 4 euros
\ei
\end{frame}

\begin{frame}{The same in a graph}
\begin{columns}[c]
\begin{column}{0.52\textwidth}
\centering
\begin{tikzpicture}[scale=0.92]
  \draw[ax] (0,0) -- (6,0) node[below left, lbl] {Income, $Y$};
  \draw[ax] (0,0) -- (0,5.2) node[above right, lbl] {Demand $Z$, production $Y$};
  \draw[guide] (0,0) -- (5.2,5.2) node[above, lbl, pelMuted] {$45^\circ$: $Y=Z$};
  \draw[curveB] (0,1.2) -- (5.6,3.44) node[right, lbl] {$ZZ$};
  \draw[shifted] (0,1.8) -- (5.6,4.04) node[right, lbl] {$ZZ'$};
  % A: Y = 1.2/0.6 = 2 ; A': 1.8/0.6 = 3
  \fill[pelText] (2,2) circle (2pt) node[left=3pt, lbl] {A};
  \fill[pelAccent] (3,3) circle (2pt) node[left=3pt, lbl] {A$'$};
  \draw[guide] (2,0) node[below, lbl] {$Y$} -- (2,2);
  \draw[guide] (3,0) node[below, lbl] {$Y'$} -- (3,3);
  \draw[{Stealth}-{Stealth}, pelAccent] (0.25,1.2) -- (0.25,1.8);
  \node[right, lbl, pelAccent] at (0.25,1.5) {$\Delta c_0$};
\end{tikzpicture}
\end{column}
\begin{column}{0.45\textwidth}
\small
\bi
  \item $ZZ$: demand as a function of income. Intercept = autonomous spending, slope = $c_1$
  \item Equilibrium where $ZZ$ crosses the $45^\circ$ line
  \item Autonomous spending rises: $ZZ$ shifts up
  \item Output rises by \textbf{more} than the shift: the horizontal gap $Y' - Y$ is larger than the vertical shift
\ei
\end{column}
\end{columns}
\end{frame}

\begin{frame}{Where the multiplier comes from}
\footnotesize
Demand rises by 1 euro. With $c_1 = 0.75$:
\vspace{0.1cm}

\centering\footnotesize
\begin{tabular}{clr}
\toprule
\textbf{Round} & \textbf{What happens} & \textbf{Extra output}\\
\midrule
1 & Firms produce 1 more euro, households earn 1 more euro & 1.00 \\
2 & Households spend $0.75$ of it & 0.75 \\
3 & \ldots of which they spend $0.75$ again & 0.56 \\
4 & & 0.42 \\
$\vdots$ & & $\vdots$ \\
\midrule
 & \textbf{Total}: $1 + c_1 + c_1^2 + \ldots = 1/(1-c_1)$ & \textbf{4.00} \\
\bottomrule
\end{tabular}
\vspace{0.15cm}

\raggedright\small
\bi
  \item Each round is smaller: the process converges
  \item It takes time: consumers and firms do not adjust instantly
  \item The size of the multiplier is an \textbf{empirical} question: we come back to it with fiscal policy (Ramey 2019)
\ei
\end{frame}

%=============================================================================
\section{Investment equals saving}
%=============================================================================

\begin{frame}{Investment equals saving}
\small
Private saving: $S \equiv Y_D - C = Y - T - C$. Public saving: $T - G$ (a budget surplus if positive, a deficit if negative)
\vspace{0.1cm}

Start from $Y = C + I + G$, subtract $T$ and $C$ from both sides:
\[
\underbrace{Y - T - C}_{S} = I + G - T
\qquad\Longrightarrow\qquad
\boxed{\,I = S + (T - G)\,}
\]
\bi
  \item Same equilibrium, seen from another side: what firms want to invest must equal what households and the government want to save
  \item This is where the name \textbf{IS} comes from: Investment = Saving
  \item With $S = -c_0 + (1-c_1)(Y-T)$ we get back exactly the same $Y$ as before
\ei
\pelTableNote{Keynes's formulation in the General Theory (1936). We ignore saving by firms (retained profits): the identity does not depend on it.}
\end{frame}

\begin{frame}{The paradox of saving}
\footnotesize
Households decide to save more at any level of income: $c_0$ falls
\bi
  \item Demand falls, so output falls (multiplier)
  \item What happens to total saving? In equilibrium $I = S + (T-G)$
  \item $\bar I$, $T$ and $G$ have not changed, so $S$ \textbf{cannot change}: income falls exactly enough to offset the wish to save more
  \item If investment also falls with output, saving ends up \emph{lower}
\ei
\vspace{0.1cm}
\begin{block}{}
Trying to save more can mean lower output and no more saving, \textbf{in the short run}
\end{block}
\bi
  \item Same for the government: cutting the deficit lowers output in the short run, unless monetary policy moves too. Over longer horizons, higher saving can raise income
\ei
\end{frame}

\begin{frame}{A warning: is the government omnipotent?}
\footnotesize
The model suggests that choosing $G$ and $T$ gives the government the output it wants. It does not, because:
\be
  \item Changing $G$ or $T$ takes time: budgets must be approved
  \item Investment, imports and the exchange rate also react, with dynamics that are hard to predict
  \item \textbf{Expectations}: a tax cut seen as temporary has a smaller effect; people may expect future taxes to pay for today's spending
  \item Pushing output too high raises inflation
  \item Deficits accumulate into public debt
\ee
\vspace{0.1cm}
\begin{block}{}
Each point is a reason why the effect of a fiscal policy must be \textbf{estimated}, not assumed: Part II is about these limits
\end{block}
\end{frame}

%=============================================================================
\section{Financial markets}
%=============================================================================

\begin{frame}{Money, income, wealth}
\small
Words with precise meanings:
\bi
  \item \textbf{Income}: what you earn in a period. A flow
  \item \textbf{Saving}: the part of disposable income you do not spend. A flow
  \item \textbf{Wealth}: the value of your assets minus your liabilities. A stock
  \item \textbf{Money}: what you can use to pay. Currency and current account deposits. A stock
\ei
\vspace{0.1cm}
Simplest financial world: only two assets
\bi
  \item \textbf{Money}: you can pay with it, but it pays no interest
  \item \textbf{Bonds}: they pay an interest rate $i$, but you cannot pay with them
\ei
How much of your wealth do you keep in money? It depends on how much you spend, and on how much interest you give up
\end{frame}

\begin{frame}{The demand for money}
\begin{columns}[c]
\begin{column}{0.5\textwidth}
\small
\[
M^d = \$Y \; L(\underset{(-)}{i})
\]
\bi
  \item Proportional to \textbf{nominal} income $\$Y$: more transactions, more money needed. Prices double, money demand doubles
  \item Decreasing in $i$: the higher the rate, the more costly it is to hold money instead of bonds
  \item Higher $\$Y$: the curve shifts right
\ei
\end{column}
\begin{column}{0.47\textwidth}
\centering
\begin{tikzpicture}[scale=0.9]
  \draw[ax] (0,0) -- (5.6,0) node[below left, lbl] {Money, $M$};
  \draw[ax] (0,0) -- (0,4.6) node[above right, lbl] {Interest rate, $i$};
  \draw[curveB] (0.6,4.2) .. controls (1.2,1.6) and (2.6,0.8) .. (4.6,0.55) node[right, lbl] {$M^d$};
  \draw[shifted] (1.5,4.2) .. controls (2.1,1.6) and (3.5,0.8) .. (5.4,0.55);
  \node[lbl, pelAccent] at (3.6,2.3) {$M^{d\prime}$ (higher $\$Y$)};
\end{tikzpicture}
\end{column}
\end{columns}
\end{frame}

\begin{frame}{Money supply and the interest rate}
\begin{columns}[c]
\begin{column}{0.5\textwidth}
\small
\bi
  \item The central bank supplies $M^s = M$
  \item Equilibrium: $M = \$Y\,L(i)$
  \item $i$ adjusts until people are willing to hold exactly the money that exists
\ei
\vspace{0.1cm}
\textbf{Comparative statics}
\bi
  \item $\$Y$ up, $M$ given: $i$ up
  \item $M$ up, $\$Y$ given: $i$ down
\ei
\end{column}
\begin{column}{0.47\textwidth}
\centering
\begin{tikzpicture}[scale=0.9]
  \draw[ax] (0,0) -- (5.6,0) node[below left, lbl] {Money, $M$};
  \draw[ax] (0,0) -- (0,4.6) node[above right, lbl] {Interest rate, $i$};
  \draw[curveB] (0.6,4.2) .. controls (1.2,1.6) and (2.6,0.8) .. (4.6,0.55) node[right, lbl] {$M^d$};
  \draw[curveA] (1.9,0) -- (1.9,4.3) node[above, lbl] {$M^s$};
  \draw[shifted] (3.0,0) -- (3.0,4.3) node[above, lbl] {$M^{s\prime}$};
  \fill[pelText] (1.9,1.57) circle (2pt) node[right, lbl] {A};
  \fill[pelAccent] (3.0,0.98) circle (2pt) node[above right, lbl] {A$'$};
  \draw[guide] (0,1.57) node[left, lbl] {$i$} -- (1.9,1.57);
  \draw[guide] (0,0.98) node[left, lbl] {$i'$} -- (3.0,0.98);
\end{tikzpicture}
\end{column}
\end{columns}
\end{frame}

\begin{frame}{How the central bank does it}
\small
\textbf{Open market operations}
\bi
  \item Central bank balance sheet: assets = bonds, liabilities = money
  \item To expand money, it \textbf{buys} bonds and pays with newly created money. To contract, it sells
\ei
\vspace{0.1cm}
\textbf{Bond prices and interest rates move in opposite directions}
\bi
  \item A one-year bond pays 100 euros in a year and costs $P_B$ today
  \[
  i = \frac{100 - P_B}{P_B} \qquad\Longleftrightarrow\qquad P_B = \frac{100}{1+i}
  \]
  \item $P_B = 98$: $i \approx 2.0\%$. $P_B = 95$: $i \approx 5.3\%$
  \item The central bank buys bonds, bond prices rise, the interest rate falls
\ei
\end{frame}

\begin{frame}{Choosing money or choosing the interest rate?}
\small
\bi
  \item We described the central bank as choosing $M$ and letting $i$ adjust
  \item Equivalently, it can choose $i$ and supply whatever $M$ is needed to get there
  \item Modern central banks do the second: the news is ``the ECB raised rates'', not ``the ECB reduced the money supply''
\ei
\vspace{0.1cm}
\textbf{One limit: the zero lower bound}
\bi
  \item When $i = 0$, money and bonds pay the same: once transactions are covered, people do not care which they hold
  \item Adding more money then does not lower $i$: a \textbf{liquidity trap}
  \item Seen as a curiosity before 2008, it became a binding constraint after the financial crisis
\ei
\pelTableNote{In practice rates went slightly below zero in several countries, including the euro area: holding large amounts of cash is costly and unsafe. The logic of the bound remains.}
\end{frame}

%=============================================================================
\section{IS-LM}
%=============================================================================

\begin{frame}{The IS curve: goods market with an interest rate}
\begin{columns}[c]
\begin{column}{0.52\textwidth}
\footnotesize
Investment now depends on sales and on the interest rate:
\[
I = I(\underset{(+)}{Y}, \underset{(-)}{i})
\]
Goods market equilibrium:
\[
Y = C(Y - T) + I(Y, i) + G
\]
\bi
  \item Higher $i$, lower investment, lower demand, lower output (multiplier)
  \item IS: the pairs $(Y, i)$ where the goods market is in equilibrium. Downward sloping
  \item $G$ up, $T$ down, confidence up: IS shifts \textbf{right}
\ei
\end{column}
\begin{column}{0.45\textwidth}
\centering
\begin{tikzpicture}[scale=0.9]
  \draw[ax] (0,0) -- (5.6,0) node[below left, lbl] {Output, $Y$};
  \draw[ax] (0,0) -- (0,4.6) node[above right, lbl] {Interest rate, $i$};
  \draw[curveB] (0.6,4.2) -- (4.4,0.5) node[right, lbl] {IS};
  \draw[shifted] (1.6,4.2) -- (5.4,0.5) node[right, lbl] {IS$'$};
  \node[lbl, pelAccent, align=left] at (4.3,3.3) {$G\uparrow$ or $T\downarrow$};
  \draw[-{Stealth}, pelAccent] (2.2,2.4) -- (3.0,2.4);
\end{tikzpicture}
\end{column}
\end{columns}
\end{frame}

\begin{frame}{The LM curve: two versions}
\small
Money market equilibrium in real terms: $\dfrac{M}{P} = Y\,L(i)$
\vspace{0.1cm}

\begin{columns}[T]
\begin{column}{0.48\textwidth}
\centering
\begin{tikzpicture}[scale=0.6]
  \draw[ax] (0,0) -- (5,0) node[below left, lbl] {$Y$};
  \draw[ax] (0,0) -- (0,4) node[above right, lbl] {$i$};
  \draw[curveA] (0.5,0.5) .. controls (2.5,1.2) and (3.6,2.2) .. (4.4,3.6) node[right, lbl] {LM};
\end{tikzpicture}

\footnotesize\raggedright
\textbf{Traditional}: the central bank fixes $M$. Higher income raises money demand, so $i$ must rise. Upward sloping
\end{column}
\begin{column}{0.48\textwidth}
\centering
\begin{tikzpicture}[scale=0.6]
  \draw[ax] (0,0) -- (5,0) node[below left, lbl] {$Y$};
  \draw[ax] (0,0) -- (0,4) node[above right, lbl] {$i$};
  \draw[curveA] (0,1.8) node[left, lbl] {$\bar i$} -- (4.6,1.8) node[right, lbl] {LM};
\end{tikzpicture}

\footnotesize\raggedright
\textbf{Blanchard, 7th ed.}: the central bank fixes $i = \bar i$ and supplies whatever money is needed. Horizontal
\end{column}
\end{columns}
\vspace{0.15cm}
\small
We use the horizontal version: closer to how the ECB and the Fed work today
\end{frame}

\begin{frame}{IS-LM: the equilibrium}
\begin{columns}[c]
\begin{column}{0.47\textwidth}
\footnotesize
\[
\begin{aligned}
\text{IS:}\quad & Y = C(Y-T) + I(Y,i) + G\\
\text{LM:}\quad & i = \bar i
\end{aligned}
\]
\bi
  \item Point A: both markets in equilibrium
  \item Output $Y_A$ is determined by demand, given the rate chosen by the central bank
\ei
\vspace{0.1cm}
\textbf{Three steps for any question}
\be
  \item Does it shift IS, LM, or both?
  \item What happens to $Y$ and $i$?
  \item Say it in words
\ee
\end{column}
\begin{column}{0.5\textwidth}
\centering
\begin{tikzpicture}[scale=0.9]
  \draw[ax] (0,0) -- (5.6,0) node[below left, lbl] {Output, $Y$};
  \draw[ax] (0,0) -- (0,4.6) node[above right, lbl] {Interest rate, $i$};
  \draw[curveB] (0.6,4.2) -- (4.6,0.3) node[right, lbl] {IS};
  \draw[curveA] (0,2) node[left, lbl] {$\bar i$} -- (5.2,2) node[right, lbl] {LM};
  % intersection: x = 0.6 + (4.2-2)*(4.0/3.9) = 2.857
  \fill[pelText] (2.857,2) circle (2pt) node[above right, lbl] {A};
  \draw[guide] (2.857,0) node[below, lbl] {$Y_A$} -- (2.857,2);
\end{tikzpicture}
\end{column}
\end{columns}
\end{frame}

\begin{frame}{Fiscal policy: a tax increase}
\begin{columns}[c]
\begin{column}{0.5\textwidth}
\small
The government raises $T$ to reduce the deficit (a \textbf{fiscal consolidation})
\be
  \item IS shifts left; LM does not move (the central bank keeps $\bar i$)
  \item The economy moves \emph{along} LM from A to A$'$: output falls
  \item Lower disposable income, lower consumption, lower output through the multiplier; lower output also cuts investment
\ee
\vspace{0.1cm}
Shift \emph{of} a curve versus movement \emph{along} a curve: keep them apart
\end{column}
\begin{column}{0.47\textwidth}
\centering
\begin{tikzpicture}[scale=0.9]
  \draw[ax] (0,0) -- (5.6,0) node[below left, lbl] {Output, $Y$};
  \draw[ax] (0,0) -- (0,4.6) node[above right, lbl] {Interest rate, $i$};
  \draw[curveB] (1.4,4.2) -- (5.4,0.3) node[right, lbl] {IS};
  \draw[shifted] (0.4,4.2) -- (4.4,0.3) node[left, lbl] {IS$'$};
  \draw[curveA] (0,2) node[left, lbl] {$\bar i$} -- (5.4,2) node[above, lbl] {LM};
  \fill[pelText] (3.657,2) circle (2pt) node[above right, lbl] {A};
  \fill[pelAccent] (2.657,2) circle (2pt) node[above left, lbl] {A$'$};
  \draw[guide] (3.657,0) node[below, lbl] {$Y$} -- (3.657,2);
  \draw[guide] (2.657,0) node[below, lbl] {$Y'$} -- (2.657,2);
\end{tikzpicture}
\end{column}
\end{columns}
\end{frame}

\begin{frame}{Monetary policy: a rate cut}
\begin{columns}[c]
\begin{column}{0.5\textwidth}
\small
The central bank lowers $\bar i$ to $\bar i'$ (a \textbf{monetary expansion}; it supplies more money to get there)
\be
  \item LM shifts down; IS does not move
  \item The economy moves \emph{along} IS from A to A$'$: output rises
  \item Lower rate, more investment, more demand and output; higher income, more consumption
\ee
\vspace{0.1cm}
Raising the rate is a \textbf{monetary contraction} (tightening): the opposite
\end{column}
\begin{column}{0.47\textwidth}
\centering
\begin{tikzpicture}[scale=0.9]
  \draw[ax] (0,0) -- (5.6,0) node[below left, lbl] {Output, $Y$};
  \draw[ax] (0,0) -- (0,4.6) node[above right, lbl] {Interest rate, $i$};
  \draw[curveB] (0.6,4.2) -- (4.6,0.3) node[right, lbl] {IS};
  \draw[curveA] (0,2.6) node[left, lbl] {$\bar i$} -- (5.2,2.6) node[right, lbl] {LM};
  \draw[shifted] (0,1.4) node[left, lbl] {$\bar i'$} -- (5.2,1.4) node[right, lbl] {LM$'$};
  % x at i=2.6: 0.6 + 1.6*4/3.9 = 2.241 ; at i=1.4: 0.6 + 2.8*4/3.9 = 3.472
  \fill[pelText] (2.241,2.6) circle (2pt) node[above right, lbl] {A};
  \fill[pelAccent] (3.472,1.4) circle (2pt) node[above right, lbl] {A$'$};
  \draw[guide] (2.241,0) node[below, lbl] {$Y$} -- (2.241,2.6);
  \draw[guide] (3.472,0) node[below, lbl] {$Y'$} -- (3.472,1.4);
\end{tikzpicture}
\end{column}
\end{columns}
\end{frame}

\begin{frame}{The policy mix}
\small
Fiscal and monetary policy used together. Why not just one?
\bi
  \item A fiscal expansion raises the deficit and, over time, the debt
  \item If $i$ is already close to zero, monetary policy has little room: fiscal policy must do the job
  \item They change different parts of output: a tax cut mostly raises $C$, a rate cut mostly raises $I$
  \item Neither works perfectly: using both reduces the risk
\ei
\vspace{0.1cm}
\centering\footnotesize
\begin{tabular}{lccccc}
\toprule
 & \textbf{IS} & \textbf{LM} & $Y$ & $C$ & $I$ \\
\midrule
Tax increase & left & -- & $\downarrow$ & $\downarrow$ & $\downarrow$ \\
Cut in $G$ & left & -- & $\downarrow$ & $\downarrow$ & $\downarrow$ \\
Rate cut & -- & down & $\uparrow$ & $\uparrow$ & $\uparrow$ \\
Rate increase & -- & up & $\downarrow$ & $\downarrow$ & $\downarrow$ \\
\bottomrule
\end{tabular}
\end{frame}

\begin{frame}{Cutting the deficit without a recession}
\begin{columns}[c]
\begin{column}{0.5\textwidth}
\small
\bi
  \item Output is about right, but the deficit $G - T$ is too large
  \item Fiscal consolidation alone: IS left, A to A$'$, a recession
  \item Add a rate cut: LM down, the economy reaches A$''$ with the \textbf{same output} as A
  \item Investment rises: same output, lower rate
\ei
\vspace{0.05cm}
\footnotesize
Blanchard's example: the US in the 1990s, Clinton (fiscal) and Greenspan (monetary). Deficit reduction raises investment \emph{only} if monetary policy accommodates it
\end{column}
\begin{column}{0.47\textwidth}
\centering
\begin{tikzpicture}[scale=0.9]
  \draw[ax] (0,0) -- (5.6,0) node[below left, lbl] {Output, $Y$};
  \draw[ax] (0,0) -- (0,4.6) node[above right, lbl] {Interest rate, $i$};
  \draw[curveB] (1.4,4.4) -- (5.4,0.5) node[right, lbl] {IS};
  \draw[shifted] (0.4,4.4) -- (4.4,0.5) node[left, lbl] {IS$'$};
  \draw[curveA] (0,2.8) node[left, lbl] {$\bar i$} -- (5.4,2.8) node[right, lbl] {LM};
  \draw[curveA, dashed] (0,1.8) node[left, lbl] {$\bar i'$} -- (5.4,1.8) node[right, lbl] {LM$'$};
  % IS at i=2.8: 1.4 + 1.6*4/3.9 = 3.041 ; IS' at 2.8: 2.041 ; IS' at 1.8: 0.4+2.6*4/3.9 = 3.067
  \fill[pelText] (3.041,2.8) circle (2pt) node[above right, lbl] {A};
  \fill[pelAccent] (2.041,2.8) circle (2pt) node[above left, lbl] {A$'$};
  \fill[pelAccent] (3.067,1.8) circle (2pt) node[below left, lbl] {A$''$};
  \draw[guide] (3.05,0) node[below, lbl] {$Y = Y''$} -- (3.05,2.8);
\end{tikzpicture}
\end{column}
\end{columns}
\end{frame}

\begin{frame}{Does IS-LM fit the facts?}
\footnotesize
The model jumps from A to A$'$. Reality takes time: consumption, investment and production all adjust with lags
\vspace{0.1cm}

Christiano, Eichenbaum and Evans (1996), US data 1960--1990: what happens after the Fed raises its rate by 1 percentage point?
\bi
  \item Retail sales fall, by about 0.9\% at the trough, after 5 quarters
  \item Output falls more slowly, by about 0.7\% after 8 quarters
  \item Employment falls, unemployment rises
  \item The price level barely moves for about 6 quarters
\ei
\vspace{0.1cm}
\begin{block}{}
Monetary policy works, with \textbf{long lags}. And prices are roughly fixed in the short run, as IS-LM assumes, but not beyond
\end{block}
\pelTableNote{As reported in Blanchard, Figure 5-10. These are estimates with confidence bands: how to identify the effect of a policy that responds to the economy is the same question we asked in Part I (Nakamura and Steinsson 2018, on the reading list).}
\end{frame}

\begin{frame}{What about aggregate demand and aggregate supply?}
\footnotesize
\bi
  \item IS-LM takes the price level as given: fine for a few quarters, not beyond
  \item Older textbooks then move to \textbf{AD-AS}: demand and supply as functions of the price level
  \item Blanchard's 7th edition drops AD-AS: he finds it cumbersome and too optimistic about how fast output returns to normal
  \item He uses \textbf{IS-LM-PC} instead: IS-LM plus a Phillips curve
  \[
  \pi - \pi_{-1} = \frac{\alpha}{L}\,(Y - Y_n)
  \]
  \bi
    \item $Y_n$: potential output. $Y - Y_n$: the output gap
    \item Output above potential: inflation rises. Below: inflation falls
  \ei
  \item Message: in the short run, demand sets output. In the medium run, output returns to potential, and monetary policy affects how fast it gets there
\ei
\pelTableNote{Blanchard, preface and chapter 9. We will use the logic, not the full model.}
\end{frame}

\begin{frame}{Why this matters for the rest of Part II}
\small
\bi
  \item In the euro area the policy rate is set by the \textbf{ECB} for all member countries: Italy cannot move its own LM
  \item A single country acts mainly through fiscal policy, its IS curve, within the \textbf{EU fiscal rules}
  \item So the questions for the next weeks:
  \bi
    \item what the EU rules allow and how they constrain $G$ and $T$ (this week)
    \item how big the multiplier is, and how we know (fiscal policy)
    \item how the ECB sets $\bar i$ and what happens at the zero lower bound (monetary policy)
  \ei
\ei
\vspace{0.1cm}
\begin{block}{}
\centering IS-LM tells us the \textbf{direction} of a policy's effect. Evaluation tells us the \textbf{size}
\end{block}
\end{frame}

\begin{frame}{To remember}
\small
\bi
  \item GDP: final goods = value added = incomes
  \item Short run: output is determined by demand. $Y = \frac{1}{1-c_1}[\,c_0 + \bar I + G - c_1 T\,]$
  \item Goods market equilibrium can also be read as $I = S + (T - G)$
  \item Money demand $M^d = \$Y\,L(i)$; the central bank sets $i$ and supplies the money needed
  \item IS: downward sloping, moved by fiscal policy. LM: horizontal at $\bar i$, moved by monetary policy
  \item Effects take time, and their size must be estimated
\ei
\vspace{0.1cm}
\textbf{Reading}: Blanchard, \textit{Macroeconomics}, 7th ed., section 2-1, chapter 3, sections 4-1, 4-2 and 4-4, chapter 5. Optional: section 9-1
\end{frame}

\begin{frame}{References}
\tiny
\bi
  \item Blanchard, O. (2017), \textit{Macroeconomics}, 7th ed., Global Edition, Pearson
  \item Christiano, L., M. Eichenbaum and C. Evans (1996), The effects of monetary policy shocks: evidence from the flow of funds, \textit{Review of Economics and Statistics}, 78(1), 16--34
  \item Nakamura, E. and J. Steinsson (2018), Identification in macroeconomics, \textit{Journal of Economic Perspectives}, 32(3), 59--86
  \item Ramey, V. A. (2019), Ten years after the financial crisis: what have we learned from the renaissance in fiscal research?, \textit{Journal of Economic Perspectives}, 33(2), 89--114
\ei
\end{frame}

\pelQuestions

\end{document}
